% Einheit 3 von 3 – Kreisteile (bank/kreis/e3.jsonl, zusammenbau v0.1)
%% TODO Zweigzeile Teil 1: was hier gelernt wird (Satz in Schülersprache aus dem Einheitstitel) – Einheit 3
\einheitenkopf[e3]{Einheit 3 von 3 · Kreisteile}
\zweigzeile{ab Kl. 7, je nach Buch bis Kl. 9 · P10}
\setcounter{aufgabe}{19}

%% TODO Titel: Ich-kann-Satz fehlt in der Bank (Platzhalter: Kettenname) – Nr. 20
%% TODO Anweisung: ein Satz über der ersten Teilaufgabe fehlt in der Bank – Nr. 20
\begin{aufgabe}{Kreisteile}
%% TODO \rechenplatz{n}: Schrittzahl des Grundfalls fehlt in der Bank (nur bei fester Schreibform, 2.3 b)
\begin{teile}
\teil Welcher Teil des Kreises ist grau? Kreuze an.\\ \kreuz{$\frac{1}{2}$}\\ \kreuz{$\frac{1}{3}$}\\ \kreuz{$\frac{1}{4}$}

\kreissektor{180}{}
\teil Ein Kreisausschnitt hat den Mittelpunktswinkel $\alpha = 90^\circ$. Welcher Teil des Kreises ist das? Gib den Anteil als Bruch an. \leerfeld
\teil Ein Kreisausschnitt hat den Mittelpunktswinkel $\alpha = 100^\circ$. Wie viel Prozent des Kreises sind das? Runde auf eine Stelle. ≈ \leerfeld[\%]
\teil Ein Kreisausschnitt ist $\frac{3}{8}$ des ganzen Kreises. Wie groß ist sein Mittelpunktswinkel? α = \leerfeld °
\teil Ein Kreisausschnitt hat den Radius $r = 9$ cm und den Mittelpunktswinkel $\alpha = 40^\circ$. Wie lang ist sein Bogen $b$? Rechne mit der $\pi$-Taste und runde auf eine Stelle. b ≈ \leerfeld[cm]
\teil Ein Kreisausschnitt hat den Radius $r = 6$ cm und den Mittelpunktswinkel $\alpha = 84^\circ$. Wie groß ist seine Fläche? Rechne mit der $\pi$-Taste und runde auf eine Stelle. A ≈ \leerfeld[cm²]
\teil Ein Kreisausschnitt hat den Radius $r = 8$ cm und den Mittelpunktswinkel $\alpha = 45^\circ$. Wie groß ist der Umfang des Ausschnitts? Rechne mit der $\pi$-Taste und runde auf eine Stelle. ≈ \leerfeld[cm]
\teil Ein Kreisring hat außen den Radius $7$ cm und innen den Radius $4$ cm. Wie groß ist seine Fläche? Rechne mit der $\pi$-Taste und runde auf eine Stelle. A ≈ \leerfeld[cm²]
\teil Im Bild ist ein Kreisausschnitt grau gefärbt. Sein Mittelpunktswinkel beträgt $130^\circ$, der Winkel des weißen Teils heißt $\alpha$. Wie viel Prozent der Kreisfläche sind grau? Runde auf eine Stelle \hfill (P10 2025 OS).

\kreissektor{130}{$130^\circ$}
≈ \leerfeld[\%]
\end{teile}
\end{aufgabe}

%% TODO Titel: Ich-kann-Satz fehlt in der Bank (Platzhalter: Kettenname) – Nr. 21
%% TODO Anweisung: ein Satz über der ersten Teilaufgabe fehlt in der Bank – Nr. 21
\begin{aufgabe}{Sachaufgabe (Rasensprenger, Tortenstück, Sektor im Kreisdiagramm)}
\begin{teile}
\teil Ein Rasensprenger spritzt $6$ m weit und dreht sich dabei hin und her über einen Winkel von $150^\circ$. Wie groß ist die Fläche, die er bewässert? Rechne mit der $\pi$-Taste und runde auf eine Stelle. ≈ \leerfeld[m²]
\end{teile}
\end{aufgabe}

%% TODO Titel: Ich-kann-Satz fehlt in der Bank (Platzhalter: Kettenname) – Nr. 22
%% TODO Anweisung: ein Satz über der ersten Teilaufgabe fehlt in der Bank – Nr. 22
\begin{aufgabe}{Kreisteile – Fehler finden, Begründen, Anwendung, Darstellungswechsel}
\begin{teile}
\teil Im Bild ist ein Ausschnitt mit $110^\circ$ grau, der Rest ist weiß. Jana berechnet den grauen Anteil: \rechnung{250^\circ : \text{Vollwinkel} &\approx 0{,}694 \\ &= 69{,}4\,\%} Wo ist der Fehler?

\kreissektor{110}{$110^\circ$}
\teil Begründe, warum ein Halbkreis zum Mittelpunktswinkel $180^\circ$ gehört.
\teil Ein Rasensprenger steht in der Ecke eines Gartens (rechter Winkel) und spritzt $7$ m weit. Auf jeden Quadratmeter sollen $4$ Liter Wasser fallen. Wie viel Wasser braucht er für einen Durchgang?
\teil Das Kreisdiagramm zeigt, wie Schüler einer Klasse zur Schule kommen. Gib zu jedem Sektor den Anteil in Prozent an.

\kreisdiagramm{Bus $162^\circ$/162, Rad $108^\circ$/108, Auto $90^\circ$/90}
\end{teile}
\end{aufgabe}
